\documentclass[11pt,letterpaper]{article}
\usepackage[T1]{fontenc}
\usepackage[utf8]{inputenc}
\usepackage{lmodern}
\usepackage[margin=1in,headheight=15pt]{geometry}
\usepackage{amsmath,amssymb,amsthm,mathtools,mathrsfs}
\usepackage{microtype,booktabs,longtable,tabularx,array,enumitem}
\usepackage{xcolor,fancyhdr,listings}
\definecolor{accent}{RGB}{63,78,56}
\definecolor{linktone}{RGB}{35,66,88}
\usepackage[colorlinks=true,linkcolor=linktone,citecolor=linktone,urlcolor=linktone]{hyperref}
\usepackage[nameinlink,noabbrev]{cleveref}
\hypersetup{pdftitle={Differential Equations over the Surcomplex Numbers},
 pdfauthor={OpenAI, prepared for Vladimir Reshetnikov},
 pdfsubject={Finite-phase obstructions, the Berarducci--Mantova derivation, and support-controlled Hahn-coherent differential equations}}
\pagestyle{fancy}\fancyhf{}
\fancyhead[L]{\small Differential equations over the surcomplex numbers}
\fancyhead[R]{\small September 2026}
\fancyfoot[C]{\thepage}
\setlength{\parskip}{0.3em}
\setlength{\parindent}{1.15em}
\setlength{\emergencystretch}{3em}
\setlist{itemsep=0.2em,topsep=0.3em,leftmargin=1.8em}
\setcounter{tocdepth}{2}
\numberwithin{equation}{section}
\newtheorem{theorem}{Theorem}[section]
\newtheorem{proposition}[theorem]{Proposition}
\newtheorem{lemma}[theorem]{Lemma}
\newtheorem{corollary}[theorem]{Corollary}
\theoremstyle{definition}
\newtheorem{definition}[theorem]{Definition}
\newtheorem{example}[theorem]{Example}
\theoremstyle{remark}
\newtheorem{remark}[theorem]{Remark}
\newcommand{\No}{\mathbf{No}}
\newcommand{\SC}{\mathbf{No}[i]}
\newcommand{\R}{\mathbb R}
\newcommand{\C}{\mathbb C}
\newcommand{\Q}{\mathbb Q}
\newcommand{\Z}{\mathbb Z}
\newcommand{\N}{\mathbb N}
\newcommand{\OO}{\mathcal O}
\newcommand{\mm}{\mathfrak m}
\newcommand{\HH}{\mathscr H}
\newcommand{\D}{\partial}
\newcommand{\Dz}{D_z}
\newcommand{\J}{\mathcal J}
\newcommand{\Obs}{\mathfrak o}
\newcommand{\Einf}{\operatorname{Exp}_{\mathrm{inf}}}
\newcommand{\Log}{\operatorname{Log}}
\newcommand{\Linf}{\operatorname{Log}_{\mathrm{inf}}}
\newcommand{\cis}{\operatorname{cis}}
\newcommand{\supp}{\operatorname{supp}}
\newcommand{\st}{\operatorname{st}}
\newcommand{\tr}{\operatorname{tr}}
\newcommand{\GL}{\operatorname{GL}}
\newcommand{\Mat}{\operatorname{Mat}}
\newcommand{\Imag}{\operatorname{Im}}
\newcommand{\Real}{\operatorname{Re}}
\newcommand{\id}{\operatorname{id}}
\newcommand{\ev}{\operatorname{ev}}
\newcommand{\ct}{\operatorname{ct}}
\newcommand{\projinf}{\pi_{\infty}}
\newcommand{\repoSHA}{\texttt{e260237db9b71da8b74a0c13c8e6355119091100}}
\newcolumntype{Y}{>{\raggedright\arraybackslash}X}
\newcolumntype{P}[1]{>{\raggedright\arraybackslash}p{#1}}
\lstset{basicstyle=\ttfamily\small,breaklines=true,columns=fullflexible,
 keepspaces=true,showstringspaces=false,frame=single,framerule=0.3pt}
\renewcommand{\UrlFont}{\small\ttfamily}

\begin{document}
\hypersetup{pageanchor=false}
\begin{titlepage}
\centering
\vspace*{0.1in}
{\large\scshape A companion to the Surreal repository\par}
\vspace{0.2in}
{\Huge\bfseries Differential Equations\par}
\vspace{0.05in}
{\Huge\bfseries over the Surcomplex Numbers\par}
\vspace{0.15in}
{\Large Finite-phase obstructions for surreal derivations\par}
\vspace{0.04in}
{\Large and support-controlled Hahn-coherent dynamics\par}
\vspace{0.15in}
{\color{accent}\rule{0.85\textwidth}{1.2pt}\par}
\vspace{0.12in}
{\large Prepared for Vladimir Reshetnikov\par}
{\large September 21, 2026\par}
\vspace{0.15in}
\begin{minipage}{0.91\textwidth}
\begin{abstract}
The repository carefully distinguishes coordinate differentiation, formal
parameter differentiation, and differentiation of surreal numbers, but its
computer-algebra discussion stops short of a solution theory connecting those
choices. This article supplies such a theory in two deliberately separate
categories.

For the Berarducci--Mantova derivation $\D$, normalized by $\D\omega=1$, we
prove an exact rank-one criterion: $\D y=(a+ib)y$ has a nonzero solution in
$\No[i]$ if and only if $b$ has a finite surreal antiderivative. Equivalently,
\[
 \left\{\frac{\D y}{y}:y\in\No[i]^{\times}\right\}
       =\No+i\D\mm.
\]
Thus algebraic closure and surjective additive integration do not supply all
complex exponential solutions. An explicit infinite-phase obstruction,
sharp power and iterated-logarithm thresholds, a harmonic-oscillator
nonexistence proof, and exact Laurent-series resolvents make the criterion
concrete.

For coordinate differentiation $D_z$, we prove existence, uniqueness, support
bounds, variation of constants, and monodromy for holomorphic matrix systems
with positive Hahn perturbations on a common ordinary domain. We also prove a
polynomially nonlinear perturbation theorem. The construction uses strong
summability, not a fictitious topological limit of surreal Picard iterates.
A noncommuting two-scale system, infinitesimal monodromy, and an infinite
pole outside a finite halo illustrate the distinctions.
\end{abstract}
\end{minipage}
\vfill
\begin{minipage}{0.91\textwidth}\small
\textbf{Status.} Established input theorems are attributed explicitly; the
connecting results are proved here. No priority claim, solution of a named
published open problem, or machine-checked formalization is asserted.
The accompanying program passes 115 finite exact symbolic checks; those
checks are not substitutes for the proofs.
\end{minipage}
\end{titlepage}
\hypersetup{pageanchor=true}
\pagenumbering{roman}
\tableofcontents
\clearpage
\pagenumbering{arabic}

\section{The documentation gap and the scope of this article}
\label{sec:audit}

\subsection{A specific gap, not a claim that differentiation is absent}
The review is pinned to the repository tree
\begin{center}\repoSHA\end{center}
retrieved on September 21, 2026. Its catalogue describes fifteen report
packages, including extensive treatments of analysis, trigonometry,
foundations, and computer algebra \cite{RepoCatalogue}.
The useful gap is \emph{not} the elementary observation that different
notions of derivative exist: the documentation already makes that point.
It is the missing development from those distinctions to precise
\emph{existence, obstruction, and support-control theorems for differential
equations}.

In particular, the computer-algebra article distinguishes the operators and
then has a short subsection titled ``Differential equations must name their
category.'' It warns about fine-local nonuniqueness and mentions formal
coefficient recursion, but does not develop a solution theory in that
subsection \cite{RepoCAS}. The trigonometry report explicitly says that it
chooses no derivation on $\No$, and that its classification of global phase
extensions is not a classification of differential-equation solutions
\cite{RepoTrig}. These are sensible boundaries. This article fills the
space between them rather than repeating the geometric material.

\begin{center}\small
\begin{tabularx}{\textwidth}{@{}P{0.22\textwidth}YY@{}}
\toprule
Inspected component & What is already present & What is developed here\\
\midrule
Computer algebra & Three derivative semantics and an initial-value warning &
A complete rank-one scalar criterion and certified exact resolvents\\[0.5em]
Trigonometry & Finite-angle polar decomposition, with no scalar derivation chosen &
Compatibility with $\D$ and a precise obstruction to infinite phase drift\\[0.5em]
Analysis catalogue and README & Common-domain coherence and strong summability &
Support-controlled linear and nonlinear initial-value theorems\\[0.5em]
Foundations catalogue & Set-sized workspaces for set-sized data &
A closure construction respecting differentiation and normalized integration\\
\bottomrule
\end{tabularx}
\end{center}

The audit is targeted: the catalogue, the analysis and trigonometry READMEs,
and the relevant computer-algebra passages were inspected. It is not a
line-by-line verification of every merged manuscript or every archived
source. Accordingly, the claim is that the identified passages leave a
substantive topic undeveloped, not that no related sentence occurs anywhere
in the repository. The included \texttt{REPOSITORY\_AUDIT.md} records the
paths, source ranges, and revision.

\subsection{The two main answers}
There are two different questions that can be hidden in the notation
$y'=ay$. In a differential field, $y$ is an element and the derivative is a
specified operation on that field. In a coherent analytic theory, $y$ is a
family of ordinary holomorphic coefficient functions and the derivative
acts on the coordinate. They have different constant fields and different
existence theorems.

For scalar surreal differentiation, the essential obstruction is not the
size of the coefficient $b$ alone. It is the size of an \emph{antiderivative}
of its imaginary part. Even $b=\omega^{-1}$ is too large in this integrated
sense, while $b=\omega^{-2}$ is permitted. For coherent coordinate equations,
positive Hahn support is a constructive resource: it orders the dependencies
of every coefficient and makes the infinite construction locally finite in
the exponent, even at higher valuation rank.

The results are exposition and deductions from known structures. The
Berarducci--Mantova construction and the $H$-field results behind it are
established mathematics \cite{BM,ADH}. The finite-angle idea is already a
central theorem of the repository \cite{RepoTrig}; we reprove the portion
needed for differential compatibility. Classical complex ODE theory is used
only for ordinary coefficient equations \cite{Teschl}. The Hahn support
argument is given in full in \cref{app:support}.

\section{Three derivatives and two kinds of constant}
\label{sec:derivatives}

\subsection{Operators must have a declared parent}
Put $t=\omega^{-1}$. The following three calculations are all correct:
\begin{equation}\label{eq:three}
 D_z(t^3)=0,\qquad D_t(t^3)=3t^2,\qquad \D(t^3)=-3t^4.
\end{equation}
In the first, $z$ is a coordinate and $t$ a coefficient. In the second, $t$
is a formal indeterminate. In the third, $t$ denotes a particular surreal
number and $\D\omega=1$, so that $\D t=-t^2$.
The notation $D_t$ becomes related to $\D$ only after specifying an embedding
and verifying that the formal operations commute with its interpretation.

Similarly, $D_z\omega=0$ but $\D\omega=1$. A number may be a coefficient
constant without being a differential-field constant. The distinction is
particularly important for initial conditions: the general solution of a
coherent linear equation is parametrized by Hahn constants, whereas the
general solution of a scalar equation for $\D$ is parametrized, when it
exists, by ordinary complex constants.

\subsection{The scalar structure used throughout}
The real surreal field $\No$ is real closed, and its complexification
$\SC$ is algebraically closed. We use the Gonshor real exponential and fix
the Berarducci--Mantova derivation. The imported properties needed here are
\begin{equation}\label{eq:BM}
\begin{gathered}
 \D\omega=1,\qquad \ker\D=\R,\qquad \D\No=\No,\\
 \D(xy)=x\D y+y\D x,\qquad
 \D\exp x=\exp x\,\D x,\\
 \D\left(\sum_j x_j\right)=\sum_j\D x_j
 \quad\text{for strongly summable set-indexed families}.
\end{gathered}
\end{equation}
The derivation is small: if $\OO$ is the convex hull of $\R$ in $\No$, then
$\D\OO\subseteq\OO$. These facts are supplied by the construction and its
Liouville-closedness theorem, not proved from the definition of surreal
numbers in this article \cite[Theorems A and B]{BM}.

Write
\[
 \OO=\{x\in\No:|x|<n\text{ for some }n\in\N\},\qquad
 \mm=\{x\in\No:|x|<1/n\text{ for every }n\ge1\}.
\]
Every finite surreal has a unique decomposition $r+\epsilon$ with
$r\in\R$ and $\epsilon\in\mm$. Complexified finite and infinitesimal
classes are $\OO+i\OO$ and $\mm+i\mm$.

\begin{lemma}[The derivative of a finite number is infinitesimal]
\label[lemma]{lem:small}
One has $\D\mm\subseteq\mm$ and $\D\OO=\D\mm\subseteq\mm$.
\end{lemma}
\begin{proof}
For positive $\epsilon\in\mm$, let $q=\sqrt\epsilon$. Then $q\in\mm$,
$\D q\in\OO$ by smallness, and
$\D\epsilon=2q\D q\in\mm$. Negative infinitesimals follow by changing
sign, and zero is immediate. Finally $\D(r+\epsilon)=\D\epsilon$ because
$\D r=0$.
\end{proof}

This elementary strengthening of smallness will be the whole obstruction
in the simplest oscillator example. It also shows why ``small derivation''
is a structural hypothesis, not an estimate about numerical rounding.

\subsection{Proper classes and set-sized constructions}
We work in a classical foundation in which the surreal class and the
specified operations can be handled as definable class objects. Every
series has set-sized support, every family being summed is set-indexed,
and every recursively constructed support below is a set. Statements about
$\No$ are class-level statements about its elements; they do not require a
set containing all surreals.

The class projection used later supplies canonical representatives of an
obstruction. We do not construct a set whose members are proper-class
cosets. \Cref{sec:workspaces} gives a direct set-sized closure argument for
any prescribed set of scalar data. This distinguishes a legitimate local
workspace from an unjustified invocation of a set-sized differential-closure
theorem on the entire proper class.

\section{Strong summability and finite analytic operations}
\label{sec:hahn}

\subsection{A valuation convention}
Let $\Gamma$ be a set-sized ordered abelian subgroup of the additive surreal
numbers, and interpret
\[
 t^\gamma=\omega^{-\gamma},\qquad
 K_\Gamma=\C((t^\Gamma)).
\]
Here $t^\gamma$ denotes a Hahn monomial; in particular this notation does
not mean that $t$ is an ordinary complex variable. An element is a formal
sum $\sum_{\gamma\in S}c_\gamma t^\gamma$ with $S$ well ordered. Its
valuation is the least exponent with nonzero coefficient; $v(0)=+\infty$.
Thus $v(t^\gamma)>0$ means that the monomial is infinitesimal. No divisibility
assumption on $\Gamma$ is needed for the ODE constructions.

A family $(F_j)_{j\in J}$ is \emph{strongly summable} if the union of its
supports is well ordered and, at every exponent, only finitely many family
members have nonzero coefficient. Summation is coefficientwise and uses
only finite sums at each exponent.

\begin{lemma}[Positive-support calculus]\label[lemma]{lem:neumann}
If $S\subseteq\Gamma_{>0}$ is well ordered, its additive monoid
\[
 S^*=\{0\}\cup S\cup(S+S)\cup(S+S+S)\cup\cdots
\]
is well ordered. For each $\gamma$, only finitely many finite words
$(s_1,\ldots,s_n)$ in $S$ have sum $\gamma$.
Consequently, if $h$ has positive support, any formal scalar power series
$\sum_{n\ge0}a_nh^n$ with ordinary coefficients is strongly summable.
\end{lemma}
A proof, including the finite-word argument, appears in \cref{app:support}.
This is the positive-support lemma underlying the classical Hahn--Neumann
construction \cite{Neumann}.

The statement is stronger than ``valuations increase.'' At higher rank,
$n\sigma$ need not eventually exceed every element of $\Gamma$. For example,
with lexicographic $\Gamma=\Z^2$, let $v(t)=(0,1)$ and $v(u)=(1,0)$.
Then $v(t^n)<v(u)$ for every ordinary $n$. The identity
$\sum_{n\ge0}t^n=(1-t)^{-1}$ is strongly summable, but its partial sums do
not converge in that valuation topology. In the full surreal fine topology
there are still more scales. Conversely, fixed rank-one Hahn fields do
have nontrivial valuation-convergent sequences. None of these topologies
should be silently identified with the others.

\subsection{Infinitesimal exponential and logarithm}
For $h\in\mm+i\mm$, define
\begin{equation}\label{eq:explog}
 \Einf(h)=\sum_{n\ge0}\frac{h^n}{n!},\qquad
 \Linf(1+h)=\sum_{n\ge1}\frac{(-1)^{n+1}h^n}{n}.
\end{equation}
Each family is strongly summable. Formal identities in finitely many
commuting variables can be substituted into positive-support Hahn series:
every coefficient computation remains finite by \cref{lem:neumann}.
In particular the two maps are inverse, exponentials multiply under
addition, and logarithms add under multiplication of elements of
$1+\mm+i\mm$.

If $h\ne0$, the leading term of $\Einf(h)-1$ is the leading term of $h$.
Indeed all powers of degree at least two have valuation strictly larger
than $v(h)$. Thus
\begin{equation}\label{eq:exp-injective}
 \Einf(h)=1\quad\Longleftrightarrow\quad h=0
 \qquad(h\in\mm+i\mm).
\end{equation}
This injectivity is local and does not assert an injective exponential on
the whole surcomplex class.

\subsection{Differentiating these sums}
The complexified derivation, constructed in the next section, is strongly
additive because its real and imaginary parts are. For infinitesimal $h$,
\begin{equation}\label{eq:infchain}
 \D\Einf(h)=\Einf(h)\D h,\qquad
 \D\Linf(1+h)=\frac{\D h}{1+h}.
\end{equation}
For example, strong additivity and the finite power rule give
\[
 \D\sum_{n\ge0}\frac{h^n}{n!}
 =\sum_{n\ge1}\frac{h^{n-1}\D h}{(n-1)!}
 =\D h\sum_{m\ge0}\frac{h^m}{m!}.
\]
The summability of the differentiated family is part of the strong-additivity
hypothesis. Multiplication by the fixed element $\D h$ is legitimate Hahn
multiplication. This argument does not apply to a Taylor series at an
arbitrary infinite argument: its original family may not be summable.

\section{Complexifying the surreal differential field}
\label{sec:complexify}

\begin{proposition}[Unique extension and its constants]\label[proposition]{prop:complexify}
The derivation on $\No$ extends uniquely to $\SC$, necessarily by
\[
 \D(x+iy)=\D x+i\D y.
\]
Its constant field is $\C$, it commutes with conjugation, and it is
surjective as an additive map.
\end{proposition}
\begin{proof}
Differentiating $i^2=-1$ gives $2i\D i=0$, hence $\D i=0$. The displayed
formula is therefore forced and directly satisfies the derivation rules.
Its kernel consists exactly of elements whose two components belong to
$\ker\D=\R$. Conjugation commutes with it by the same formula. Given
$a+ib$, choose real surreal antiderivatives of $a$ and $b$ and combine them.
\end{proof}

\begin{proposition}[Logarithmic derivative and modulus]\label[proposition]{prop:modulus}
For $y\ne0$, put $r=|y|=\sqrt{y\overline y}>0$. Then
\begin{equation}\label{eq:radial}
 \frac{\D r}{r}=\Real\frac{\D y}{y}.
\end{equation}
If $\D y=(a+ib)y$, then $\D r=ar$ and $u=y/r$ satisfies
$|u|=1$ and $\D u=ibu$.
\end{proposition}
\begin{proof}
Differentiate $r^2=y\overline y$ and divide by $2r^2$. This gives
$\D r/r=\tfrac12(\D y/y+\overline{\D y/y})$.
The equation for $u$ follows from the quotient rule.
\end{proof}

Additive surjectivity solves $\D y=f$ for every $f\in\SC$. It does
\emph{not} solve $\D y=fy$ for every $f$ with $y\ne0$. In a real
Liouville-closed exponential field the latter equation has a positive
solution for every real $f$, namely $\exp A$ with $\D A=f$.
Complexification introduces imaginary coefficients, and their effect is
not determined by real exponentiation.

\section{Finite phases and their differential compatibility}
\label{sec:phases}

\subsection{Constructing the finite phase}
For a finite real surreal $\theta=r+\epsilon$, where $r\in\R$ and
$\epsilon\in\mm$, define
\begin{equation}\label{eq:cis}
 \cis\theta=e^{ir}\Einf(i\epsilon).
\end{equation}
The ordinary value $e^{ir}$ belongs to $\C$ and is a differential constant.
The real and imaginary parts of \eqref{eq:cis} are the finite surreal cosine
and sine. This agrees with their Taylor construction.

\begin{theorem}[Finite-angle polar decomposition]\label{thm:polar}
The map
\[
 \cis:(\OO,+)\longrightarrow
       \{u\in\SC:u\overline u=1\}
\]
is a surjective homomorphism with kernel $2\pi\Z$. Every nonzero surcomplex
number therefore has a representation $y=r\cis\theta$, with $r>0$ and
$\theta$ finite. Moreover,
\begin{equation}\label{eq:cisderivative}
 \D\cis\theta=i(\D\theta)\cis\theta.
\end{equation}
\end{theorem}
\begin{proof}
The homomorphism property follows from ordinary complex exponentiation on
the real parts and the infinitesimal exponential identity. Conjugation
changes the sign of the angle, so the image has unit modulus.

For surjectivity, let $u\overline u=1$. Its real and imaginary parts have
absolute value at most one and hence have standard parts. Put
$u_0=\st(u)\in\C$; then $|u_0|=1$, and
$u/u_0=1+h$ with $h\in\mm+i\mm$. Let
$L=\Linf(1+h)$. Since $\overline{1+h}=(1+h)^{-1}$, the formal logarithm
identities imply $\overline L=-L$. Thus $L=i\delta$ with $\delta\in\mm$.
Choose an ordinary real $r$ with $e^{ir}=u_0$. Then
$u=e^{ir}\Einf(i\delta)=\cis(r+\delta)$.

If $\cis(r+\epsilon)=1$, taking standard parts gives $r\in2\pi\Z$.
Then \eqref{eq:exp-injective} gives $\epsilon=0$. This proves the kernel
statement. Polar decomposition follows by applying surjectivity to $y/|y|$.
Finally \eqref{eq:infchain} and $\D r=0$ give
\eqref{eq:cisderivative}.
\end{proof}

The surjectivity and kernel statement are already central to the repository's
trigonometry \cite{RepoTrig}. The role here is different: finite angles
supply a differential constraint on \emph{every} surcomplex unit, even those
obtained by some noncanonical global phase convention.

\subsection{Why an infinite phase cannot be hidden in a unit}
A unit $u$ may be denoted by an expression involving a large parameter.
Nevertheless, as an element of $\SC$ it has a finite angle $\theta$. Hence
\[
 \frac{\D u}{u}=i\D\theta\in i\D\mm.
\]
The set $\D\mm$ is contained in the infinitesimals by \cref{lem:small}, but
that containment is not an equality. A coefficient such as $\omega^{-1}$
is infinitesimal while its primitives are infinite. A notation like
$e^{i\log\omega}$ cannot, by itself, turn that coefficient into a permissible
logarithmic derivative.

This is not a restriction on ordinary complex exponentiation, where the
variable is a coordinate. Nor is it a claim that arbitrary global phase
characters are impossible. It is a restriction on compatibility with the
\emph{specified scalar derivation}.

\section{The exact rank-one solvability criterion}
\label{sec:rankone}

\subsection{A normalized antiderivative}
Every surreal has a Conway normal form. Let $\ct(x)$ denote the ordinary
real coefficient of the monomial $1=\omega^0$; this is defined even when
$x$ is infinite. It is an $\R$-linear projection. Let $\projinf(x)$ be the
subseries consisting of monomials strictly larger than $1$.

\begin{definition}[Normalized primitive and infinite-phase obstruction]
\label[definition]{def:obstruction}
For $b\in\No$, let $\J b$ be the unique surreal satisfying
\[
 \D\J b=b,\qquad \ct(\J b)=0.
\]
Define $\Obs(b)=\projinf(\J b)$.
\end{definition}

Existence follows by taking any primitive and subtracting its constant
coefficient. Uniqueness follows because two primitives differ by a real
constant. Thus $\J$ and $\Obs$ are $\R$-linear class functions defined
without choosing representatives from a proper-class collection of cosets.
The definition is exact but is \emph{not} a claim of a total effective
antidifferentiation algorithm.

\begin{theorem}[Rank-one phase criterion]\label{thm:rankone}
Let $a,b\in\No$. The following are equivalent:
\begin{enumerate}
\item There is a nonzero $y\in\SC$ satisfying $\D y=(a+ib)y$.
\item The element $b$ has a finite real surreal antiderivative.
\item $b\in\D\mm$.
\item $\Obs(b)=0$.
\end{enumerate}
When these conditions hold, all solutions are
\begin{equation}\label{eq:all-solutions}
 y=c\,\exp(\J a)\Einf(i\J b),\qquad c\in\C.
\end{equation}
The nonzero solutions correspond to $c\in\C^\times$.
\end{theorem}
\begin{proof}
Suppose first that $y\ne0$ solves the equation. By \cref{thm:polar}, write
$y=r\cis\theta$ with $r>0$ and $\theta$ finite. Differentiation gives
\[
 \frac{\D y}{y}=\frac{\D r}{r}+i\D\theta.
\]
Its imaginary part is $b$, so $\D\theta=b$. This proves (1)$\Rightarrow$(2).
A finite primitive differs from an infinitesimal primitive by its real
standard part, proving equivalence of (2) and (3).
Because $\ct(\J b)=0$, the normal-form decomposition says that
$\J b$ is finite exactly when it is infinitesimal, and exactly when
$\projinf(\J b)=0$. This gives equivalence with (4).

Under these conditions, $\J b\in\mm$. Thus the right side of
\eqref{eq:all-solutions} is defined. The real exponential chain rule and
\eqref{eq:infchain} show that the factor with $c=1$ is a nonzero solution.
This proves the remaining implication. If $y_0$ is that solution and $y$ is
any other one, the quotient rule gives $\D(y/y_0)=0$. By
\cref{prop:complexify}, $y/y_0\in\C$. Zero is included by $c=0$.
\end{proof}

\begin{corollary}[The image of the logarithmic derivative]
\label[corollary]{cor:logimage}
The group homomorphism $y\mapsto\D y/y$ has kernel $\C^\times$ and image
\begin{equation}\label{eq:logimage}
 \No+i\D\mm.
\end{equation}
In particular, it is not onto $\SC$.
\end{corollary}
\begin{proof}
The product rule makes it a homomorphism from multiplication to addition.
Its kernel is the nonzero constant field. The image follows from
\cref{thm:rankone}; it is proper since $1\notin\D\mm$ by
\cref{lem:small}.
\end{proof}

\subsection{The obstruction is a complete additive invariant}
Let $\Pi$ denote the class of purely infinite surreal numbers, including
zero. The map $\Obs:\No\to\Pi$ is onto: for $p\in\Pi$,
$\J(\D p)=p$ and hence $\Obs(\D p)=p$. Its kernel is exactly $\D\mm$.
Thus the class-level sequence
\begin{equation}\label{eq:exact}
 0\longrightarrow\D\mm\longrightarrow\No
     \xrightarrow{\ \Obs\ }\Pi\longrightarrow0
\end{equation}
is exact. This is shorthand for the assertions about the displayed maps
and their elements; no set of all classes is required.

The normal-form splitting of $\J b$ also gives the unique decomposition
\[
 b=\D\projinf(\J b)+\D\bigl(\J b-\projinf(\J b)\bigr),
\]
where the second summand lies in $\D\mm$.
The obstruction therefore identifies precisely the part of an imaginary
coefficient requiring an infinite phase primitive. It is more informative
than the necessary inequality $b\in\mm$.

\subsection{Inhomogeneous equations}
\begin{proposition}[Variation of constants when the phase is admissible]
\label[proposition]{prop:inhom}
Suppose $\Obs(b)=0$, let $y_0$ be the nonzero solution in
\eqref{eq:all-solutions} with $c=1$, and let $f\in\SC$. Then
\[
 \D y=(a+ib)y+f
\]
has all solutions
\[
 y=y_0\bigl(\J_{\C}(f/y_0)+c\bigr),\qquad c\in\C,
\]
where $\J_{\C}(r+is)=\J r+i\J s$.
\end{proposition}
\begin{proof}
The product rule verifies the formula, and any two solutions differ by a
solution of the homogeneous equation.
\end{proof}

When $\Obs(b)\ne0$, the homogeneous kernel is zero. An inhomogeneous
solution, if it exists, is therefore unique, but the obstruction to a
\emph{nonzero homogeneous} solution does not itself prohibit one. For
example, $y=1$ solves $\D y=iy-i$. No general inhomogeneous-surjectivity
claim on the whole class is needed here; \cref{sec:resolvents} supplies an
explicit surjectivity theorem on a useful Laurent workspace.

\section{Sharp thresholds and the obstruction to oscillation}
\label{sec:examples}

\subsection{A power-law threshold}
All powers in this subsection have \emph{ordinary real} exponents and are
interpreted by the real exponential and logarithm. No constant-exponent
rule for an arbitrary surreal exponent is being used.

\begin{proposition}[Power threshold]\label[proposition]{prop:power}
For $p\in\R$, the equation
\[
 \D y=i\omega^{-p}y
\]
has a nonzero surcomplex solution if and only if $p>1$.
\end{proposition}
\begin{proof}
The chain rule gives the primitives
\[
 B_p=\begin{cases}
 \displaystyle\frac{\omega^{1-p}}{1-p},&p\ne1,\\[0.6em]
 \log\omega,&p=1.
 \end{cases}
\]
For $p>1$ the first expression is infinitesimal. For $p<1$ its magnitude is
infinite, and $\log\omega$ is infinite as well. All other primitives differ
by an ordinary real constant, so their finiteness is unchanged. Apply
\cref{thm:rankone}.
\end{proof}

The threshold is stricter than infinitesimality of the coefficient. The
following three equations have different outcomes:
\begin{align*}
 \D y=iy&\quad\text{has only the zero solution},\\
 \D y=i\omega^{-1}y&\quad\text{also has only the zero solution},\\
 \D y=i\omega^{-2}y&\quad\text{has }y=c\Einf(-i\omega^{-1}),\quad c\in\C.
\end{align*}
The first obstruction is visible from smallness alone. The second needs the
integrated criterion: $\omega^{-1}$ is small but its primitive is not.
Adding real growth does not remove an imaginary obstruction. For example,
\[
 \D y=(1+i\omega^{-2})y
 \quad\Longrightarrow\quad
 y=c\exp(\omega)\Einf(-i\omega^{-1}),
\]
whereas replacing $\omega^{-2}$ by $\omega^{-1}$ leaves no nonzero solution.

\subsection{Successive logarithmic boundaries}
Let $L_0=\omega$ and $L_{j+1}=\log L_j$, for ordinary finite $j$.
Every $L_j$ is positive infinite. The identities
\[
 \D L_k=\frac{1}{L_0\cdots L_{k-1}}
 \quad(k\ge1)
\]
follow inductively from the logarithm chain rule.

\begin{proposition}[Iterated-logarithm threshold]\label[proposition]{prop:logthreshold}
For an ordinary integer $k\ge0$ and $p\in\R$, put
\[
 b_{k,p}=\frac{1}{L_0L_1\cdots L_{k-1}L_k^p},
\]
where the empty product is $1$. The equation $\D y=ib_{k,p}y$ has a
nonzero surcomplex solution exactly when $p>1$.
\end{proposition}
\begin{proof}
A primitive is $L_k^{1-p}/(1-p)$ for $p\ne1$, and $L_{k+1}$ for $p=1$.
The first is finite precisely when $p>1$, while the second is infinite.
Again all primitives differ only by a real constant.
\end{proof}

Thus passing to a smaller coefficient by inserting an additional logarithm
creates a new boundary rather than eliminating the phenomenon. These are
finite iterated constructions; no infinite product or transfinite logarithm
tower is required by the proposition.

\subsection{A direct harmonic-oscillator proof}
\begin{theorem}[No nonzero scalar harmonic oscillator]\label{thm:oscillator}
In $\No$, and consequently in $\SC$, the equation
\[
 \D^2 f+f=0
\]
has only the zero solution.
\end{theorem}
\begin{proof}
First suppose $f$ is real. The energy
$E=f^2+(\D f)^2$ satisfies
\[
 \D E=2\D f\,(f+\D^2f)=0,
\]
so $E\in\R_{\ge0}$. It follows that $f$ and $\D f$ are finite.
By \cref{lem:small}, $\D f$ and $\D^2 f$ are infinitesimal. Since
$f=-\D^2 f$, the element $f$ is infinitesimal too. Hence $E$ is both
real and infinitesimal, and is therefore zero. Positivity gives $f=0$.
For complex $f$, apply this argument to its real and imaginary parts.
\end{proof}

The argument isolates a structural difference from ordinary real calculus.
A real-valued function can remain bounded while its ordinary derivative
oscillates. Here a \emph{finite field element} has infinitesimal scalar
derivative. Energy conservation and smallness together prohibit that
oscillation. This does not say that coordinate functions $\sin z$ and
$\cos z$ fail their ordinary differential equations.

\subsection{Why algebraic closure is not differential closure}
The complex field $\SC$ is algebraically closed, and its derivation is
additively surjective. Nevertheless $\D y=iy$, with $y\ne0$, is not
solvable in it. These facts concern different closure properties.
An explicit differential extension is easy on a set-sized base: start with
$\C(\omega)$, adjoin a transcendental indeterminate $E$, and prescribe
\[
 \D\omega=1,\qquad \D E=iE.
\]
The polynomial product rule and then the quotient rule define a derivation
on $\C(\omega)(E)$. There is no differential embedding of this extension
into $\SC$ fixing $\C(\omega)$, because the image of $E$ would be a
forbidden nonzero solution. We need not make any assertion about the full
constant field of the constructed extension to obtain this conclusion.

\begin{corollary}[No unrestricted scalar-compatible complex exponential]
\label[corollary]{cor:noexp}
There is no map $\mathcal E:\SC\to\SC^\times$ satisfying
\[
 \D\mathcal E(z)=\mathcal E(z)\D z\qquad\text{for every }z\in\SC.
\]
\end{corollary}
\begin{proof}
Set $z=i\omega$. Its image would be a nonzero solution of $\D y=iy$.
\end{proof}

No continuity, group law, or normalization was needed in this corollary.
It is therefore not repaired by a clever choice of infinite circular phase.
It also does not contradict a global exponential with a \emph{coordinate}
chain rule, or a group-theoretic extension of finite trigonometry. Those are
precisely the distinctions the repository's derivative warning requires.

\section{Set-sized workspaces closed under differential operations}
\label{sec:workspaces}

The full surreal class is convenient for statements but unnecessary for a
prescribed set of scalar calculations. A workspace must, however, be closed
under the operations that will actually be used. Merely naming an exponent
group does not guarantee closure under the Berarducci--Mantova derivation.

\begin{theorem}[Differential and integration closure for set-sized data]
\label{thm:workspace}
Given a set $A\subseteq\SC$, there is a set-sized real-closed subfield
$F\subseteq\No$ such that $A\subseteq F[i]$, $\R\subseteq F$, and $F$ is
closed under $\D$, $\J$, the real exponential, the logarithm on positive
elements, and the real and imaginary parts of $\cis\theta$ for finite
$\theta\in F$. In particular,
\[
 \ker(\D|_F)=\R,\qquad \D F=F,
 \qquad \ker(\D|_{F[i]})=\C.
\]
The rank-one criterion of \cref{thm:rankone} holds internally in $F[i]$:
for $a,b\in F$, a nonzero solution exists in $F[i]$ if and only if $b$
has a finite primitive in $\No$.
\end{theorem}
\begin{proof}
Let $F_0$ be the real closure inside $\No$ of the field generated by
$\R$ and the real and imaginary parts of $A$. Given $F_n$, form the field
generated by $F_n$ and all its images under the stated operations, using
only positive arguments for logarithms and finite arguments for $\cis$.
Let $F_{n+1}$ be its real closure inside $\No$.

Each stage is a set. Images under the specified definable class functions
are sets by replacement. Field generation uses finite expressions, and
algebraic closure inside a field uses only the roots of a set of
polynomials, each of finite degree. Hence
$F=\bigcup_{n<\omega}F_n$ is a set-sized field closed under the operations.
It is real closed because every finite polynomial calculation occurs at
some stage and its real-algebraic witnesses occur at the next. The assertions
about constants follow from the ambient constant fields. Since $\J F\subseteq F$
and $\D\J f=f$, the derivation on $F$ is onto.

If a nonzero solution exists internally, necessity follows from the ambient
criterion. Conversely, if $b$ has a finite ambient primitive, then $\J b$
is infinitesimal and belongs to $F$. Closure gives
$\exp(\J a)\cis(\J b)\in F[i]$, which is the required nonzero solution.
\end{proof}

This is a closure construction, not an effective enumeration algorithm:
normalization of primitives need not be computable from a supplied symbolic
representation. Nor does the theorem say that $F$ contains every Hahn series
over its own monomials. An arbitrary-support closure is a different demand.
The much stronger structural connections with transseries and universal
$H$-fields belong to the literature \cite{ADH}; they are not prerequisites
for this elementary set-sized construction.

A useful warning comes from the Laurent field $\R((t))$ with
$t=\omega^{-1}$. It is stable under $\D=-t^2D_t$, but $t$ has no primitive
there. Its ambient primitive $\log\omega$ requires enlargement. Algebraic,
differential, and integration closure are separate properties even in this
very concrete example.

\section{Exact Laurent resolvents and factorial solutions}
\label{sec:resolvents}

\subsection{A rank-one valuation workspace}
In this section alone let
\[
 L=\C((t)),\qquad t=\omega^{-1},\qquad \D=-t^2D_t.
\]
The last equality is valid on this embedded Laurent field by strong
additivity and the finite integer-power rule. Thus
\begin{equation}\label{eq:laurentD}
 \D\left(\sum_{n\ge m}c_nt^n\right)
       =-\sum_{n\ge m}n c_nt^{n+1}.
\end{equation}
In particular,
\begin{equation}\label{eq:valD}
 v(\D f)\ge v(f)+1
\end{equation}
for nonzero $f$, allowing $v(0)=+\infty$ on the left.
This is a special rank-one estimate, not a general valuation formula for
arbitrary surreal monomials.

\begin{proposition}[Integration in the Laurent parent]\label[proposition]{prop:laurentint}
The image $\D L$ consists exactly of Laurent series with zero coefficient
at $t^1$. For real infinitesimal Laurent series,
\[
 \D\bigl(t\R[[t]]\bigr)=t^2\R[[t]].
\]
\end{proposition}
\begin{proof}
In \eqref{eq:laurentD}, the coefficient of $t^1$ comes from $n=0$ and is
zero. Conversely, if $f=\sum c_kt^k$ has $c_1=0$, then
\[
 -\sum_{k\ne1}\frac{c_k}{k-1}t^{k-1}
\]
is a Laurent primitive. Restricting to powers $k\ge2$ proves the second
claim.
\end{proof}

\subsection{Inverting any constant-coefficient operator with nonzero constant term}
\begin{theorem}[Strong Laurent resolvent]\label{thm:resolvent}
Let $P(s)\in\C[s]$ with $P(0)\ne0$, and write the formal reciprocal
\[
 \frac1{P(s)}=\sum_{n\ge0}q_ns^n\in\C[[s]].
\]
For every $f\in L$, the family $(q_n\D^n f)_{n\ge0}$ is strongly summable
and
\begin{equation}\label{eq:resolvent}
 P(\D)^{-1}f=\sum_{n\ge0}q_n\D^nf.
\end{equation}
Consequently $P(\D):L\to L$ is bijective.
\end{theorem}
\begin{proof}
For $f\ne0$, repeated application of \eqref{eq:valD} gives
$v(\D^nf)\ge v(f)+n$. The union of supports is bounded below in $\Z$,
and a fixed exponent can occur for only finitely many $n$. Hence the family
is strongly summable. Strong additivity permits applying each of the
finitely many derivatives in $P(\D)$ termwise. At every exponent the
calculation is finite, so the formal identity $P(s)\sum q_ns^n=1$ proves
$P(\D)y=f$.

For injectivity, let $y\ne0$. The leading term of $P(\D)y$ is $P(0)$
times the leading term of $y$: every positive-order derivative has strictly
larger valuation. It cannot vanish.
\end{proof}

The theorem includes both $P(s)=s-i$ and $P(s)=s^2+1$. Thus the absence of a
nonzero homogeneous oscillator is compatible with unique solvability of
every inhomogeneous oscillator equation whose forcing lies in $L$.
There is no conflict with \cref{prop:laurentint}: that proposition concerns
$P(s)=s$, whose constant term is zero.

If $Q_N(s)=\sum_{n=0}^Nq_ns^n$ and $y_N=Q_N(\D)f$, then
$P(s)Q_N(s)-1$ has order at least $N+1$. Hence
\begin{equation}\label{eq:residual}
 v\bigl(P(\D)y_N-f\bigr)\ge v(f)+N+1.
\end{equation}
This is a checkable residual certificate for a finite output. It does not
assert that $y_N$ equals the full solution.

\subsection{A factorial series that is an exact surcomplex solution}
For $P(s)=s-i$,
\[
 \frac1{s-i}=i\sum_{n\ge0}(-i)^ns^n.
\]
Since $\D^nt=(-1)^nn!t^{n+1}$, the unique solution in $L$ of
$\D y-iy=t$ is
\begin{equation}\label{eq:factorial}
 y=\sum_{n\ge0}i^{n+1}n!t^{n+1}
   =it-t^2-2it^3+6t^4+24it^5-120t^6-\cdots.
\end{equation}
The ordinary complex power series obtained by replacing $t$ by a complex
variable has radius of convergence zero. Nevertheless its displayed family
is strongly summable at the surreal $t=\omega^{-1}$, and
\eqref{eq:factorial} is an exact field element satisfying the equation.
No choice of analytic summation method is involved in this assertion.

This example tests two independent boundaries. A coefficient growth
condition that is essential for an ordinary analytic germ is unnecessary
for this Hahn value. Conversely, interpreting the field element as a
solution does not provide an ordinary analytic germ at $t=0$.
The program supplied with this article checks finite jets and their
residuals; the infinite claim follows from \cref{thm:resolvent}.

\section{Common-domain Hahn-coherent linear systems}
\label{sec:coherent}

We now change categories. Throughout this section $D_z$ acts on an ordinary
coordinate and leaves every Hahn coefficient constant. No scalar derivation
on $K_\Gamma$ is involved.

\subsection{The ring and its evaluation}
Let $U\subseteq\C$ be a connected ordinary open set. Define
\[
 \HH_\Gamma(U)=\OO(U)((t^\Gamma))
 =\left\{\sum_{\gamma\in S}f_\gamma(z)t^\gamma:
 S\text{ well ordered},\ f_\gamma\in\OO(U)\right\}.
\]
Every coefficient is holomorphic on the \emph{same} $U$. This is not the
larger ring of unrelated convergent germs with shrinking radii, and it is
not the full formal ring $\C[[z]]((t^\Gamma))$. The distinction is already
emphasized in the repository \cite{RepoCatalogue,RepoAnalysis}.
For vector and matrix series, support means the union of the supports of
the finitely many entries. Differentiation is coefficientwise:
\[
 D_zF=\sum_\gamma f'_\gamma(z)t^\gamma.
\]
A nonzero coefficient function means one that is not identically zero.
Thus global support is not redefined each time a coefficient happens to
vanish at a particular $z$.

Evaluation at an ordinary $z_0\in U$ is a ring homomorphism to $K_\Gamma$,
and
\begin{equation}\label{eq:coherentconstants}
 \ker\bigl(D_z:\HH_\Gamma(U)\to\HH_\Gamma(U)\bigr)=K_\Gamma.
\end{equation}
Indeed an ordinary holomorphic function with zero derivative on connected
$U$ is constant. If $U$ is simply connected, define the normalized primitive
coefficientwise by
\[
 I_{z_0}F=\sum_\gamma
 \left(\int_{z_0}^{z}f_\gamma(\zeta)\,d\zeta\right)t^\gamma.
\]
The ordinary integrals are path independent and holomorphic. They do not
enlarge the support.

For completeness, the section also has a value at a finite-halo point
$z_0+h$ with $h$ infinitesimal, after enlarging the exponent workspace if
necessary:
\[
 F^\sharp(z_0+h)=
 \sum_{\gamma,n\ge0}\frac{f_\gamma^{(n)}(z_0)}{n!}t^\gamma h^n.
\]
The union of supports lies in $S+(\supp h)^*$. The two-set sum lemma and the
finite-word assertion in \cref{app:support} give strong summability of the
double family. This is an infinitesimal Taylor evaluation, not evaluation
at an arbitrary infinite coordinate.

\subsection{The linear initial-value theorem}
\begin{theorem}[Positive-perturbation fundamental matrix]
\label{thm:coherentlinear}
Let $U$ be simply connected, $z_0\in U$, and
\[
 A=A_0+E\in\Mat_d\HH_\Gamma(U),
 \qquad A_0\in\Mat_d\OO(U),
 \qquad S=\supp E\subseteq\Gamma_{>0}.
\]
There is a unique Hahn-coherent matrix $X$ satisfying
\begin{equation}\label{eq:fundamental}
 D_zX=AX,\qquad X(z_0)=I_d.
\end{equation}
It is invertible in $\Mat_d\HH_\Gamma(U)$, its support is contained in
$S^*$, and its coefficient at exponent $0$ is the ordinary normalized
fundamental matrix $\Phi_0$ of $D_z\Phi_0=A_0\Phi_0$.
Uniqueness holds among all Hahn-coherent matrices, not only those whose
support was prescribed in advance.
\end{theorem}
\begin{proof}
Ordinary holomorphic linear ODE theory gives $\Phi_0$ on $U$, normalized at
$z_0$, with everywhere nonzero determinant. Local holomorphic existence and
uniqueness, analytic continuation along ordinary paths, and simple
connectedness give the global matrix \cite[Chapter 4]{Teschl}.
Put
\[
 B=\Phi_0^{-1}E\Phi_0,
 \qquad TH=I_{z_0}(BH).
\]
The support of $B$ is contained in $S$. For each $n$, the support of
$T^nI_d$ lies in the sums of $n$ members of $S$. For a fixed exponent,
only finitely many words in $S$ contribute, by \cref{lem:neumann}. Hence
\begin{equation}\label{eq:peano}
 V=\sum_{n\ge0}T^nI_d
\end{equation}
is strongly summable, has support in $S^*$, and satisfies
$V=I_d+TV$. Consequently $D_zV=BV$ and $V(z_0)=I_d$.
Then $X=\Phi_0V$ satisfies \eqref{eq:fundamental}.

Write $V=I_d+W$, where $W$ has positive support. The geometric matrix
series $\sum_{n\ge0}(-W)^n$ is strongly summable by the same support lemma
and is a two-sided inverse. Thus $X$ is invertible and has the claimed
constant coefficient.

For uniqueness, let $Z$ be the difference of two solutions and put
$H=\Phi_0^{-1}Z$. If $H\ne0$, choose the least exponent $\gamma$ in its
global support. The equation $D_zH=BH$ says that its coefficient satisfies
$H'_\gamma=0$: any contribution on the right would require a coefficient
of $H$ at $\gamma-s<\gamma$ for some $s>0$. The initial condition gives
$H_\gamma(z_0)=0$, so $H_\gamma=0$, a contradiction.
\end{proof}

The familiar iterated-integral form in \eqref{eq:peano} does not import an
ordinary convergence proof into the surreal fine topology. The primitive
is an ordinary coefficientwise primitive; the infinite sum is a strongly
summable Hahn family. At higher rank, the successive support levels need
not become cofinal in the value group.

The proof gives an explicit recursion. Write $V_0=I_d$ and let all
unlisted coefficients be zero. For $\gamma>0$,
\begin{equation}\label{eq:linear-recursion}
 V'_\gamma=\sum_{s+\delta=\gamma}B_sV_\delta,
 \qquad V_\gamma(z_0)=0.
\end{equation}
All dependencies satisfy $\delta<\gamma$, and every displayed sum is finite.
This is a support certificate as well as an existence proof.

\subsection{Variation of constants and determinant}
\begin{corollary}[Inhomogeneous coherent systems]\label[corollary]{cor:coherentinhom}
Under the hypotheses of \cref{thm:coherentlinear}, for any
$f\in\HH_\Gamma(U)^d$ and $c\in K_\Gamma^d$, the unique solution of
$D_zy=Ay+f$, $y(z_0)=c$, is
\[
 y=X\left(c+I_{z_0}(X^{-1}f)\right).
\]
Its support is contained in
\[
 \bigl(\supp c+S^*\bigr)\ \cup\
 \bigl(\supp f+S^*\bigr).
\]
\end{corollary}
\begin{proof}
The ordinary product rule holds coefficientwise because each Hahn
coefficient is a finite sum. It verifies the formula. The support follows
from the support of $X$ and $X^{-1}$ and from $S^*+S^*=S^*$.
Uniqueness follows from the homogeneous theorem applied to differences.
\end{proof}

\begin{corollary}[Liouville determinant formula]\label[corollary]{cor:determinant}
With the same normalization,
\[
 \det X=\det\Phi_0\,
 \Einf\!\left(I_{z_0}(\tr E)\right).
\]
Here $\Einf$ is the formal exponential of a positive-support coherent
section, not an unrestricted scalar complex exponential.
\end{corollary}
\begin{proof}
Finite matrix algebra gives Jacobi's formula
$D_z\det X=(\tr A)\det X$. The proposed expression satisfies the same
scalar equation and has value $1$ at $z_0$; uniqueness in
\cref{thm:coherentlinear} with $d=1$ proves the identity. The exponent has
positive support, so its exponential is admitted by \cref{lem:neumann}.
\end{proof}

\subsection{A sharp scalar obstruction, but not a matrix necessity claim}
\begin{proposition}[A negative leading exponent forbids a coherent scalar exponential]
\label[proposition]{prop:negativecoherent}
If $a\in\HH_\Gamma(U)$ has $v(a)<0$, then $D_zY=aY$ has no nonzero
solution in $\HH_\Gamma(U)$.
\end{proposition}
\begin{proof}
Differentiation can delete coefficient functions but cannot create smaller
exponents, so $v(D_zY)\ge v(Y)$. Since $U$ is connected, $\OO(U)$ is an
integral domain, and therefore $v(aY)=v(a)+v(Y)<v(Y)$ for $Y\ne0$.
The equation is impossible.
\end{proof}

In particular $D_zY=i\omega Y$ has no nonzero common-domain Hahn-coherent
solution on any connected ordinary domain. A sufficiently microscopic
formal germ is a different matter; see \cref{sec:bridge}.
The positive-support hypothesis is \emph{not} a necessary condition for an
arbitrary matrix system. If $N^2=0$, then
\[
 D_zX=t^{-1}NX,\qquad X(0)=I_d,
 \qquad X=I_d+zt^{-1}N
\]
is a coherent polynomial matrix despite its negative exponent. Nilpotence
allows finite algebraic cancellation that the scalar argument excludes.

\section{A polynomially nonlinear perturbation theorem}
\label{sec:nonlinear}

A nonlinear theorem needs both a domain on which the ordinary background
solution exists and a reason why every new Hahn coefficient is determined
by earlier ones. Positivity supplies the latter. Polynomial dependence on
the dependent variables makes all compositions unambiguous without an
additional multivariable analytic evaluation theory.

\begin{theorem}[Support-controlled nonlinear deformation]
\label{thm:nonlinear}
Let $U\subseteq\C$ be simply connected, $z_0\in U$, and suppose
\[
 f_0(z,Y)\in\OO(U)[Y_1,\ldots,Y_d]^d,
 \qquad y_0\in\OO(U)^d,
 \qquad y'_0=f_0(z,y_0).
\]
Let $S\subseteq\Gamma_{>0}$ be well ordered and, for each $s\in S$, let
$f_s\in\OO(U)[Y_1,\ldots,Y_d]^d$. Consider
\begin{equation}\label{eq:nonlinear}
 D_zy=f_0(z,y)+\sum_{s\in S}t^sf_s(z,y),
 \qquad y(z_0)=y_0(z_0)+\epsilon,
\end{equation}
where $\epsilon\in K_\Gamma^d$ has positive support $T$.
Put $P=S\cup T$. There is a unique solution $y=y_0+u$ with $u$ a
positive-support vector in $\HH_\Gamma(U)^d$. Its support satisfies
\[
 \supp u\subseteq P^*\setminus\{0\}.
\]
No common bound on the polynomial degrees of the $f_s$ is required.
Uniqueness is among solutions with the specified ordinary reduction $y_0$.
\end{theorem}
\begin{proof}
Put
\[
 A(z)=\left.\frac{\partial f_0}{\partial Y}\right|_{Y=y_0(z)}.
\]
For any proposed positive-support $u$, polynomial expansion gives
\[
 f_0(z,y_0+u)=f_0(z,y_0)+Au+R_0(z,u),
\]
where every monomial of $R_0$ has degree at least two in $u$.
The perturbation $t^sf_s(z,y_0+u)$ carries the additional positive
exponent $s$.

Construct $u_\gamma$ by transfinite recursion on the well-ordered set
$P^*\setminus\{0\}$. At exponent $\gamma$, the equation has the form
\begin{equation}\label{eq:nonlinear-coefficient}
 u'_\gamma=A(z)u_\gamma+R_\gamma(z),
 \qquad u_\gamma(z_0)=\epsilon_\gamma.
\end{equation}
Here $R_\gamma$ depends only on coefficients $u_\delta$ with
$0<\delta<\gamma$. Indeed a contribution from $R_0$ has at least two
positive exponents whose sum is $\gamma$, so each is smaller; a contribution
from a perturbation also has $s>0$, so every remaining exponent is smaller.
The only term containing $u_\gamma$ is the displayed ordinary linear
term $Au_\gamma$.

At a fixed $\gamma$, only finitely many contributions occur. One way to see
this is to expand each previously admitted exponent as a finite word in
$P$: all contributions then determine words in $P$ with sum $\gamma$,
and \cref{lem:neumann} permits only finitely many such words. There are
only finitely many vector indices at each fixed length. Each individual
$f_s$ has finitely many polynomial monomials, and only finitely many
$s$ can participate in these words. Thus unbounded degrees across different
$s$ do not introduce an infinite coefficient sum.

By induction, $R_\gamma$ is an ordinary holomorphic vector on $U$.
The ordinary linear inhomogeneous initial-value problem
\eqref{eq:nonlinear-coefficient} has a unique holomorphic solution on $U$,
by variation of constants with the ordinary fundamental matrix of $A$.
This completes the recursion. The resulting support is contained in $P^*$,
so it is well ordered. Coefficientwise verification of
\eqref{eq:nonlinear-coefficient} proves \eqref{eq:nonlinear}.

For uniqueness, two positive-support solutions have a well-ordered union of
supports. At the least exponent of their nonzero difference, all nonlinear
and positive perturbation contributions to the difference vanish; the
remaining coefficient solves $w'=Aw$ with $w(z_0)=0$. Ordinary uniqueness
forces that coefficient to be zero, a contradiction.
\end{proof}

The theorem does not assert existence of an ordinary background solution on
an arbitrary prescribed domain. That solution is part of the hypotheses;
it may encounter poles or other obstructions beyond $U$. Nor does the
theorem cover unrestricted substitution into an arbitrary holomorphic
function of $Y$ without specifying its domain and support conditions.
These restrictions make its exact content usable rather than merely formal.

\subsection{A coherent entire solution with a pole at an infinite coordinate}
Take $d=1$, $f_0(Y)=Y^2$, $y_0=0$, no coefficient perturbation, and initial
value $\epsilon=t$. Then
\begin{equation}\label{eq:riccati}
 y=\frac{t}{1-tz}=\sum_{n\ge0}t^{n+1}z^n
\end{equation}
satisfies $D_zy=y^2$, $y(0)=t$. Every coefficient is a polynomial in $z$, so
this is a coherent section on the whole ordinary complex plane. Direct
formal differentiation verifies the equation, or it follows from
\cref{thm:nonlinear}.

The rational expression has a pole at the infinite point $z=t^{-1}=\omega$.
That point is not in the finite halo of $\C$. Evaluating the Hahn series
there would produce infinitely many copies of $t$, and is not a strongly
summable evaluation. Thus being coherent on the whole ordinary plane does
not mean being regular at every surcomplex coordinate. The location of the
pole explains, rather than contradicts, the exact geometric-series identity.

\section{Noncommuting systems and infinitesimal monodromy}
\label{sec:monodromy}

\subsection{A genuinely two-scale, noncommuting example}
Let $\Gamma=\Z^2$ with lexicographic order. Use the monomials
\[
 v(t)=(0,1),\qquad v(u)=(1,0),\qquad u\prec t^n
 \quad\text{for every ordinary }n\ge1.
\]
This is a Hahn workspace for coordinate differentiation: no formula for
$\D u$ is assumed. Set
\[
 N=\begin{pmatrix}0&1\\0&0\end{pmatrix},\qquad
 M=\begin{pmatrix}0&0\\1&0\end{pmatrix},\qquad
 A(z)=tN+uzM.
\]
Both matrices are nilpotent, but $NM\ne MN$. The normalized fundamental
matrix has the expansion
\[
 X(z)=\sum_{p,q\ge0}t^pu^qX_{pq}(z),\qquad X_{00}=I_2,
\]
with the recursion
\begin{equation}\label{eq:twoscale}
 X'_{pq}=N X_{p-1,q}+zM X_{p,q-1},\qquad
 X_{pq}(0)=0\quad(p+q>0),
\end{equation}
where a negative index means zero. The first terms are
\[
 X_{10}=zN,\qquad X_{01}=\frac{z^2}{2}M,
 \qquad
 X_{11}=z^3\left(\frac{NM}{6}+\frac{MN}{3}\right)
       =\begin{pmatrix}z^3/6&0\\0&z^3/3\end{pmatrix}.
\]
These values follow by ordinary polynomial integration of the recursion.
Since $\tr A=0$, \cref{cor:determinant} gives $\det X=1$ exactly.

It would be incorrect to replace $X$ by
$\Einf(tzN+uz^2M/2)$. The latter's mixed coefficient is
$z^3(NM+MN)/4$, not $X_{11}$. Their difference at $tu$ is
\[
 \frac{z^3}{12}(MN-NM)\ne0.
\]
The order of multiplication inside iterated integrals matters even though
every individual coefficient calculation is finite.

The verification program computes all $X_{pq}$ with $p+q\le6$, verifies the
coefficient equations and initial conditions, and checks $\det X=1$ to the
same total degree. \emph{Total degree is not a valuation cutoff here.}
For example, $t^{100}$ has higher total degree but lower valuation than $u$.
The support theorem justifies the complete series; the finite jet is a
different, explicitly declared representation of part of it.

\subsection{Monodromy remains an ordinary-domain construction}
On a nonsimply connected ordinary domain $U$, lift the coefficient system
to its ordinary universal covering surface and choose a lift of the
basepoint. The proof of \cref{thm:coherentlinear} applies there: ordinary
holomorphic primitives of coefficient one-forms are path independent on
the cover, and the same support calculation is unchanged.

Analytic continuation of the normalized fundamental matrix around an
ordinary based loop $\ell$ yields another fundamental matrix $X_\ell$.
Their ratio has zero coordinate derivative, so
\begin{equation}\label{eq:monodromy}
 X_\ell=X M_\ell,\qquad M_\ell\in\GL_d(K_\Gamma).
\end{equation}
The coefficient at exponent zero is the ordinary monodromy matrix of
$A_0$, and all other exponents lie in $S^*\setminus\{0\}$.
Indeed $M_\ell$ can be read off at the lifted endpoint, since the original
normalization is $X(z_0)=I_d$. Under a fixed convention for composition of
based loops these matrices form the usual monodromy representation; reversing
the loop-composition convention reverses the multiplication order.

No fine-continuous surcomplex path is asserted. The loop is an ordinary
complex path, continuation occurs coefficientwise, and its result is a
matrix of Hahn constants. This is the same kind of domain separation that
makes coefficientwise contour theory meaningful.

\subsection{A residue invisible to ordinary monodromy}
Consider, on a branch domain for the ordinary logarithm containing $1$,
\begin{equation}\label{eq:scalar-monodromy}
 D_zY=\frac{\lambda+\epsilon}{z}Y,\qquad Y(1)=1,
 \qquad \lambda\in\C,\quad v(\epsilon)>0.
\end{equation}
The coherent solution is
\[
 Y=z^\lambda\Einf(\epsilon\Log z),
 \qquad z^\lambda=\exp(\lambda\Log z)
 \quad\text{in ordinary complex analysis}.
\]
A positive loop about zero changes $\Log z$ by $2\pi i$ and gives
\begin{equation}\label{eq:infmonodromy}
 M=e^{2\pi i\lambda}\Einf(2\pi i\epsilon).
\end{equation}
If $\lambda\in\Z$, ordinary monodromy is $1$, but the full monodromy is
not $1$ when $\epsilon\ne0$. By \eqref{eq:exp-injective}, its leading
change is exactly $2\pi i\epsilon$.
Thus reduction to the ordinary coefficient can erase a genuine global
obstruction to a single-valued coherent solution.

Here $\epsilon$ is constant for $D_z$, not necessarily constant for the
scalar $\D$. The finite-angle obstruction of \cref{thm:rankone} is
therefore not the criterion for solving \eqref{eq:scalar-monodromy}.
For instance $\epsilon=\omega^{-1}$ is permitted in this coordinate
problem although it is prohibited as the imaginary coefficient of the
scalar equation $\D y=i\omega^{-1}y$.

\section{Formal initial conditions and the total chain rule}
\label{sec:bridge}

\subsection{Formal existence is a third, weaker assertion}
Let $K$ be any characteristic-zero field. For
$A(Z)=\sum_{n\ge0}A_nZ^n\in\Mat_d(K[[Z]])$, the formal problem
\[
 D_ZX=A(Z)X,\qquad X(0)=I_d
\]
has a unique solution $X=\sum_{n\ge0}X_nZ^n$, given by
\begin{equation}\label{eq:formal-recursion}
 X_0=I_d,\qquad
 (n+1)X_{n+1}=\sum_{k=0}^n A_kX_{n-k}.
\end{equation}
This is a direct recursion: division by $n+1$ is permitted, and the right
side uses only known coefficients. No topology or analytic estimate is
needed. It follows that the formal problem with $K=\SC$ is always
solvable in this sense.

The result does not assert that the solution can be evaluated at a given
surcomplex number. For $A=i\omega$, it is
\[
 X(Z)=\sum_{n\ge0}\frac{(i\omega Z)^n}{n!}.
\]
At any nonzero ordinary $z$, the valuations of its terms are unbounded
below, so this displayed family is not strongly summable. At an
infinitesimal $h$ with $\omega h$ infinitesimal, it is strongly summable
and equals $\Einf(i\omega h)$. Thus formal solvability, microscopic
evaluability, and common-domain coherence are three different promises.
The negative-exponent argument in \cref{prop:negativecoherent} excludes
not just this particular series representation, but any nonzero coherent
solution in the stated ring.

\subsection{Evaluating a coordinate changes the derivative calculation}
For a polynomial with surcomplex coefficients,
\[
 F(Z)=\sum_{n=0}^N a_nZ^n,
 \qquad (\D_{\mathrm{coeff}}F)(Z)=\sum_{n=0}^N(\D a_n)Z^n,
\]
one has the exact total chain rule
\begin{equation}\label{eq:totalchain}
 \D\bigl(F(x)\bigr)
   =(\D_{\mathrm{coeff}}F)(x)+(\D x)(D_ZF)(x).
\end{equation}
It follows immediately by applying the product rule to each $a_nx^n$.
For example, $F(Z)=tZ$ and $x=\omega$ give
\[
 \D(F(\omega))=(-t^2)\omega+t=0,
\]
as they must, since $F(\omega)=1$. Keeping only the coordinate term would
return the incorrect answer $t$.

For a strongly summable series, a corresponding chain rule requires that
the original evaluation and the differentiated families be admissible, and
that the coefficient derivation be defined in the parent. The polynomial
identity does not waive these obligations. On the specific ring
$\OO(U)((t))$ one may define
\[
 \D_{\mathrm{coeff}}=-t^2D_t,\qquad D_z
\]
coefficientwise, and they commute because both sides have identical
coefficient formulas. This statement need not hold as a meaningful
endomorphism assertion on an arbitrarily chosen $K_\Gamma$: the scalar
Berarducci--Mantova derivative can require monomials outside that workspace.

\subsection{A comparison that prevents false transports}
\begin{center}\small
\begin{tabularx}{\textwidth}{@{}P{0.19\textwidth}YYY@{}}
\toprule
Category & Constants & Constructive positive result & Boundary\\
\midrule
Scalar $(\SC,\D)$ & $\C$ & Rank-one criterion and finite-phase formula &
$\D y=iy$ has no nonzero solution\\[0.5em]
Laurent scalar $L$ & $\C$ & $P(\D)$ invertible when $P(0)\ne0$ &
$t$ has no Laurent primitive\\[0.5em]
Formal $K[[Z]]$ & $K$ & Linear coefficient recursion &
Evaluation may not be summable\\[0.5em]
Common-domain $\HH_\Gamma(U)$ & $K_\Gamma$ & Positive-perturbation linear
and polynomially nonlinear theorems & Negative scalar leading exponent
prevents a coherent exponential\\
\bottomrule
\end{tabularx}
\end{center}

These are not competing answers to one problem. They are answers to
precisely typed problems that happen to share differential notation.

\section{Computational consequences and limits of the results}
\label{sec:computation}

\subsection{What an exact solver can honestly return}
A solver should retain the coefficient parent, the derivative operator,
the support certificate, and the solution category as part of its result.
For a scalar rank-one equation, a positive answer can carry real primitives
$A,B$ with $\D A=a$, $\D B=b$ and a proof that $B$ is finite. A negative
answer can carry the nonzero infinite part of a normalized primitive of
$b$, or a special-case obstruction such as \cref{prop:power}. An unevaluated
primitive or an unknown finiteness test is not a negative solvability proof.

On the Laurent parent, $P(0)\ne0$ is a finite algebraic certificate. The
output can be an exact descriptor \eqref{eq:resolvent}, accompanied by a
requested finite truncation and residual bound \eqref{eq:residual}.
Factorial coefficient growth is acceptable in that exact descriptor; it
would not be acceptable as a claim of an ordinary analytic Taylor germ.

For coherent systems, a suitable result consists of an ordinary background
fundamental matrix or background nonlinear solution, a declared common
domain, a positive-support set or effective support descriptor, and the
coefficient recursion. In an effective finite-generator support monoid,
requested coefficients can often be computed by enumerating their finite
dependencies. The mathematical proof that the dependencies are finite does
not by itself give a terminating enumeration algorithm for an arbitrary
program-described ordered group. That distinction is already important in
the repository's computer-algebra capability taxonomy \cite{RepoCAS}.

A matrix jet specified by total degree, a Laurent jet specified by valuation,
and an exact recursively defined series must remain different output types.
Likewise, branch continuation and monodromy should use ordinary-domain path
data rather than an implied fine-continuous path in the surreal class.

\subsection{What has and has not been established}
The scalar theory in this article completely characterizes rank-one
\emph{homogeneous} multiplicative equations for the fixed derivation. It
also gives variation of constants in the admissible case and a full class
of inhomogeneous constant-coefficient resolvents in a Laurent field.
It does not classify arbitrary nonlinear scalar differential equations or
all higher-order operators on the full surcomplex class.

The coherent theory proves ordinary-domain linear existence, uniqueness,
invertibility, variation of constants, determinant identities, and
monodromy under a positive perturbation hypothesis. Its nonlinear theorem
requires polynomial dependence and a supplied ordinary background solution.
It does not establish a general irregular-singularity classification,
a Stokes classification, or a Riemann--Hilbert correspondence over arbitrary
Hahn fields. Matrix systems with negative exponents need finer analysis:
the nilpotent example demonstrates why the scalar obstruction is insufficient.

The set-sized closure result is mathematical, not an implemented symbolic
backend. The example program neither constructs arbitrary surreals nor
formalizes the Berarducci--Mantova construction. No mathematical novelty
claim rests on the fact that a topic was missing from the inspected
repository passages.

\subsection{Concrete next research targets}
Several extensions can now be stated without conflating categories.
For scalar differential algebra, a useful next step is to characterize
higher-rank systems by invariant data generalizing the finite-phase
obstruction; eigenvalues alone cannot encode a noncommuting variable
coefficient system. For coherent equations, negative-exponent systems can
be separated into finite nilpotent cases, cases reducible by a gauge
transformation, and genuinely noncoherent cases. Each proposed reduction
must carry its own support proof.

For computation, the immediate target is narrower and testable: implement
the Laurent resolvent and a finite-generator version of the coefficient
recursions, returning certificates in the formats above. For formalization,
one can first work with a set-sized Hahn field and the finite-word lemma,
then prove the coherent theorems independently of any full-class surreal
construction. The scalar phase theorem can be formalized abstractly from
an ordered exponential differential field with the stated smallness,
constant-field, and infinitesimal-series properties, before specializing
to the Berarducci--Mantova model.

\appendix
\section{A complete positive-support argument}
\label{app:support}

This appendix supplies the combinatorial support facts used throughout.
It works in a set-sized ordered abelian group and uses no completeness
assumption on the group. The finite-word principle is the well-ordered
alphabet case of Higman's lemma \cite{Higman}; a proof is included to make
the Hahn constructions independently checkable.

\subsection{Two elementary facts about well-ordered supports}
First, every sequence in a well-ordered set has an infinite nondecreasing
subsequence. To see this, choose an occurrence of the least value in the
whole tail; then choose an occurrence of the least value in the tail
strictly after that index, and continue. The selected indices increase and
the selected values do not decrease.

\begin{lemma}[Sum of two supports]\label[lemma]{lem:twosupports}
If $A,B$ are well-ordered subsets of an ordered abelian group, then $A+B$
is well ordered. For each $\gamma$, only finitely many pairs
$(a,b)\in A\times B$ satisfy $a+b=\gamma$. A finite union of
well-ordered subsets is well ordered.
\end{lemma}
\begin{proof}
If $a_n+b_n$ were strictly decreasing, select a nondecreasing subsequence
of the $a_n$. The corresponding $b_n$ would be strictly decreasing,
contradicting well-ordering of $B$. In a linear order, absence of an
infinite decreasing sequence implies well-ordering in the classical
set-theoretic setting used here: a nonempty subset with no least element
would produce such a sequence by recursive choice.

If infinitely many distinct pairs had a fixed sum, select a sequence of
distinct pairs. Their $a$ entries are distinct, so a nondecreasing
subsequence is strictly increasing; its $b$ entries strictly decrease.
Finally, an infinite decreasing sequence in a finite union has an infinite
subsequence in one member of the union, which is impossible.
\end{proof}

\subsection{Finite words over a well-ordered alphabet}
A word $a_1\cdots a_m$ \emph{embeds} in a word $b_1\cdots b_n$ if there
are indices $1\le j_1<\cdots<j_m\le n$ with $a_k\le b_{j_k}$ for every
$k$. The empty word embeds in every word.

\begin{lemma}[Finite-word lemma]\label[lemma]{lem:higman}
Every infinite sequence of finite words over a well-ordered alphabet has
two words, in their original order, such that the earlier embeds in the
later.
\end{lemma}
\begin{proof}
Call a sequence \emph{bad} if it has no such pair. Suppose a bad sequence
exists. Choose a bad sequence $w_0,w_1,\ldots$ minimally in the following
sense: after a fixed prefix $w_0,\ldots,w_{n-1}$, choose $w_n$ with
smallest possible length among words that permit continuation to an
infinite bad sequence. Natural-number lengths make each choice possible.
No $w_n$ is empty.

Write $w_n=v_na_n$, where $a_n$ is its last letter. Select indices
$n_0<n_1<\cdots$ for which the $a_{n_j}$ are nondecreasing. Consider
\[
 w_0,\ldots,w_{n_0-1},v_{n_0},v_{n_1},v_{n_2},\ldots.
\]
This sequence cannot be bad, since $v_{n_0}$ is shorter than the minimal
choice $w_{n_0}$ after the same prefix. A good pair in the displayed
sequence cannot lie wholly in the prefix, because the original sequence
was bad. If it goes from a prefix word to a $v_{n_j}$, it also embeds in
$w_{n_j}$, another contradiction. Finally, if it goes from $v_{n_i}$ to
$v_{n_j}$ with $i<j$, append the nondecreasing last letters to the embedding:
$w_{n_i}=v_{n_i}a_{n_i}$ embeds in $v_{n_j}a_{n_j}=w_{n_j}$.
This is again impossible. Thus no bad sequence exists.
\end{proof}

\subsection{From words to Hahn sums}
Now let $S\subseteq\Gamma_{>0}$ be well ordered. If a word in $S$ embeds
in another, the sum of its letters is at most the sum of the other's.
Moreover equality forces the words to be identical. Any unmatched letter
is strictly positive; and any strict comparison of matched letters
contributes a strictly positive difference. Their finite sum cannot vanish.

An infinite strictly decreasing sequence of word sums would therefore be
a bad word sequence, contradicting \cref{lem:higman}. Hence $S^*$ is
well ordered. If infinitely many distinct words had the same sum, select
an infinite sequence of them. An embedded pair would force identical words,
a contradiction. Thus each exponent has only finitely many word
representations. This proves \cref{lem:neumann}.

For clarity, this statement counts \emph{ordered} words, not merely
unordered multisets. It therefore suffices even for noncommuting matrix
products: each word can retain the order in which its coefficient matrices
were multiplied. Finite matrix dimensions only introduce finitely many
intermediate indices for a word of fixed length.

Combining this result with \cref{lem:twosupports} justifies multiplication
of an arbitrary Hahn series by a positive-support power series, the
Taylor-evaluation family in \cref{sec:coherent}, and every coefficient
recursion used in the article. It does not establish cofinal growth of
$n\sigma$ in a higher-rank value group, and no such assertion was used.

\section{Reproducibility and the exact scope of the checks}
\label{app:verification}

The package contains \texttt{verify\_examples.py} and its recorded
\texttt{verification\_report.json}. Running
\begin{lstlisting}
python verify_examples.py
\end{lstlisting}
with Python and SymPy produces the report. The delivered run used Python
3.13.5 and SymPy 1.14.0 and passed all 115 checks:
\begin{center}\small
\begin{tabularx}{0.94\textwidth}{@{}Yr@{}}
\toprule
Check group & Passed\\
\midrule
Power primitives & 7\\
Finite-phase formula identities & 3\\
Iterated-logarithm primitives & 15\\
Laurent resolvent residuals & 10\\
Noncommuting matrix coefficient equations and initial values & 56\\
Matrix ODE and determinant jets & 5\\
Nonlinear rational solution and jets & 15\\
Infinitesimal monodromy jets & 2\\
Total chain-rule identities & 2\\
\midrule
Total & 115\\
\bottomrule
\end{tabularx}
\end{center}

The Laurent tests use reciprocal-series order $12$, the matrix tests use
total degree $6$ in the two formal parameters, and the nonlinear test uses
powers through $t^{12}$. All algebra is symbolic and exact; there are no
floating-point samples masquerading as identities.

The primitive tests replace $\omega$ by an ordinary positive variable and
check the finite algebraic differential identities. They do not infer
that an arbitrary ordinary exponential is a scalar-compatible surcomplex
exponential. The phase-admissibility and nonexistence statements are proved
in the text, not checked by SymPy. Likewise, the program does not verify
Higman's lemma, arbitrary transfinite recursion, or the external
Berarducci--Mantova theorem.

The LaTeX source has an embedded bibliography and no external graphics or
bibliography database. Build it with
\begin{lstlisting}
latexmk -pdf -interaction=nonstopmode -halt-on-error article.tex
\end{lstlisting}
The accompanying README records the package contents and build requirements.
The repository audit is pinned rather than asserted to describe future
versions of the documentation.

\clearpage
\phantomsection
\addcontentsline{toc}{section}{References}
\begin{thebibliography}{99}
\raggedright

\bibitem{BM}
Alessandro Berarducci and Vincenzo Mantova,
\emph{Surreal numbers, derivations and transseries},
Journal of the European Mathematical Society \textbf{20} (2018), 339--390.
\href{https://doi.org/10.4171/JEMS/769}{doi:10.4171/JEMS/769}.
Preprint: \href{https://arxiv.org/abs/1503.00315}{arXiv:1503.00315}.
Theorem A supplies the differential structure; Theorem B supplies the
normalized small, surjective derivation used here.

\bibitem{ADH}
Matthias Aschenbrenner, Lou van den Dries, and Joris van der Hoeven,
\emph{The surreal numbers as a universal $H$-field},
Journal of the European Mathematical Society \textbf{21} (2019), 1179--1199.
\href{https://doi.org/10.4171/JEMS/858}{doi:10.4171/JEMS/858}.
Preprint: \href{https://arxiv.org/abs/1512.02267}{arXiv:1512.02267}.

\bibitem{Teschl}
Gerald Teschl,
\emph{Ordinary Differential Equations and Dynamical Systems},
Graduate Studies in Mathematics, vol.~140, American Mathematical Society,
2012. Chapter 4 concerns differential equations in the complex domain.
\href{https://www.mat.univie.ac.at/~gerald/ftp/book-ode/ode.pdf}
{Author's authorized preliminary manuscript}.

\bibitem{Neumann}
B.~H. Neumann,
\emph{On ordered division rings},
Transactions of the American Mathematical Society \textbf{66} (1949),
202--252.
\href{https://doi.org/10.1090/S0002-9947-1949-0032593-5}
{doi:10.1090/S0002-9947-1949-0032593-5}.
The support lemma used in this article is proved in Appendix~A.

\bibitem{Higman}
Graham Higman,
\emph{Ordering by divisibility in abstract algebras},
Proceedings of the London Mathematical Society (3) \textbf{2} (1952),
326--336.
\href{https://doi.org/10.1112/plms/s3-2.1.326}
{doi:10.1112/plms/s3-2.1.326}.

\bibitem{RepoCatalogue}
Surreal repository (maintainer: Vladimir Reshetnikov),
\emph{Documentation catalogue},
\path{docs/README.md}, audited at revision
\texttt{e260237db9b71da8b74a0c13c8e6355119091100}, September 21, 2026.
\href{https://github.com/VladimirReshetnikov/Surreal/blob/e260237db9b71da8b74a0c13c8e6355119091100/docs/README.md}
{Pinned repository source}.

\bibitem{RepoCAS}
Surreal repository (maintainer: Vladimir Reshetnikov),
\emph{Surreal and Surcomplex Numbers in Symbolic
Computer Algebra}, AI-assisted merged research report,
\path{docs/foundations-and-computation/computer-algebra/article.tex},
same audited revision. Relevant inspected source ranges: 3100--3680 and
3580--3820; in particular the derivative discussion and the subsection
``Differential equations must name their category.''
\href{https://github.com/VladimirReshetnikov/Surreal/blob/e260237db9b71da8b74a0c13c8e6355119091100/docs/foundations-and-computation/computer-algebra/article.tex}
{Pinned repository source}.

\bibitem{RepoTrig}
Surreal repository (maintainer: Vladimir Reshetnikov),
\emph{Trigonometry on the Surcomplex Plane:
Canonical Phases, Arbitrary-Scale Geometry, and Infinitesimal Degeneration},
AI-assisted merged research report, September 21, 2026.
The audit used \path{docs/surcomplex/trigonometry/README.md}, same revision.
\href{https://github.com/VladimirReshetnikov/Surreal/blob/e260237db9b71da8b74a0c13c8e6355119091100/docs/surcomplex/trigonometry/README.md}
{Pinned report description and explicit non-claims}.

\bibitem{RepoAnalysis}
Surreal repository (maintainer: Vladimir Reshetnikov),
\emph{Surcomplex analysis: holomorphic functions on
$K=\No[i]$}, AI-assisted merged research report, September 21, 2026.
The audit used \path{docs/surcomplex/analysis/README.md}, same revision.
\href{https://github.com/VladimirReshetnikov/Surreal/blob/e260237db9b71da8b74a0c13c8e6355119091100/docs/surcomplex/analysis/README.md}
{Pinned report description and function-class distinctions}.

\end{thebibliography}
\end{document}
